\renewcommand{\bibname}{Quellenverzeichnis}
% \bibliographystyle{baalphadin}

\pdfbookmark[-1]{Quellenverzeichnis}{pdf-Bibliography}%
\phantomsection%
\addcontentsline{toc}{chapter}{Quellenverzeichnis}%
% BibTeX:
% \bibliography{bibliography}\label{Bibliography}
% BibLateX:
\printbibliography

\clearpage\section*{Hilfsmittel}\label{Tools}
Zur Erstellung der vorliegenden Abschlussarbeit wurden die folgenden Hilfsmittel verwendet:
% Seien Sie so spezifisch wie möglich!
\begin{itemize}
	\item ChatGPT --\,\url{https://chat.openai.com/chat},
    \item CircuiTikZ (Version 1.7.1) --\,\url{https://ctan.dcc.uchile.cl/graphics/pgf/contrib/circuitikz/doc/circuitikzmanual.pdf},
	\item Gemeinsame Bibliothek der BA Dresden und der ehs Dresden,
	\item Google-Scholar --\,\url{https://scholar.google.com},
	\item Microsoft Windows 10 Enterprise (22H2 Build 19045.5371, x64),
	\item Overleaf --\,\url{https://www.overleaf.com},
    \item Perplexity --\,\url{https://www.perplexity.ai/},
    \item Wikibooks \LaTeX --\,\url{https://en.wikibooks.org/wiki/LaTeX}.
\end{itemize}
